\section{\texorpdfstring{\(\sigma\)-代数把可观察事件组织成封闭系统}{sigma-代数把可观察事件组织成封闭系统}}
\label{sec:12-Real-Analysis-Support/05-Measure-Integration-and-Conditional-Expectation/01}

线上服务的日志按分钟聚合：每分钟一行，记下请求总数与错误总数，别的不留。运维问，第 37 秒那一瞬的错误率是多少。

这个问题没有答案。不是概率太小，不是采样不够，是这份日志里根本没有能回答它的东西。把粒度改成每秒一行，问题立刻有了答案；再问第 37.5 秒，又没有了。粒度画到哪里，可回答的问题就到哪里为止。

同样的现象在连续概率里以更硬的形式出现。想在 \([0,1]\) 上定义一个``均匀长度''，要求平移不变，且可数多个互不相交的集合，长度之和等于并集的长度。直觉说这不难：总长度 1 的线段，均匀摊开就是了。Vitali 在 1905 年证明它做不到——他构造出一族子集，无论怎样给它们分配数值，平移不变与可数可加必有一条被破坏。这不是构造技巧不够精巧，是逻辑上不可能。实数轴上存在一些集合，谈论它们的长度这件事本身没有一致的说法。

于是连续概率模型只能退一步：先声明哪些集合可以被询问，再在这些集合上赋值。声明的载体就是 \(\sigma\)-代数。

\subsection{可回答的问题，是信息块拼出来的}

先在有限空间里把这件事看清楚。设 \(\Omega=\set{\omega_1,\omega_2,\omega_3,\omega_4}\)，测量设备粗糙，只能告诉你结果落在
\[
B_1=\set{\omega_1,\omega_2}\quad\text{还是}\quad B_2=\set{\omega_3,\omega_4}.
\]
此时可以回答的问题只有四个：什么都没发生、落在 \(B_1\)、落在 \(B_2\)、总之发生了。写成集合族就是
\[
\mathcal G=\set{\varnothing,\;B_1,\;B_2,\;\Omega}.
\]
``结果是 \(\omega_1\)''不在这张清单上。\(\omega_1\) 当然可能发生，但现有信息无法把它与 \(\omega_2\) 分开，所以``是 \(\omega_1\) 吗''这句话在当前框架下接不住。

\begin{center}
\begin{tikzpicture}[>=stealth]
\begin{scope}
  \fill[gray!30] (0,0) rectangle (1,2.2);
  \fill[gray!30] (2,0) rectangle (3,2.2);
  \draw (0,0) rectangle (4,2.2);
  \draw[very thick] (2,0) -- (2,2.2);
  \node at (1,-0.45) {\footnotesize \(B_1\)};
  \node at (3,-0.45) {\footnotesize \(B_2\)};
  \node at (2,2.65) {\footnotesize 粗粒度：两块};
  \node at (2,-1.15) {\footnotesize 灰色区域切开了两块，问不出来};
\end{scope}
\begin{scope}[shift={(6,0)}]
  \fill[gray!30] (0,0) rectangle (1,2.2);
  \fill[gray!30] (2,0) rectangle (3,2.2);
  \draw (0,0) rectangle (4,2.2);
  \draw[very thick] (1,0) -- (1,2.2);
  \draw[very thick] (2,0) -- (2,2.2);
  \draw[very thick] (3,0) -- (3,2.2);
  \node at (0.5,-0.45) {\footnotesize \(C_1\)};
  \node at (1.5,-0.45) {\footnotesize \(C_2\)};
  \node at (2.5,-0.45) {\footnotesize \(C_3\)};
  \node at (3.5,-0.45) {\footnotesize \(C_4\)};
  \node at (2,2.65) {\footnotesize 细粒度：四块};
  \node at (2,-1.15) {\footnotesize 灰色区域等于 \(C_1\cup C_3\)，可测};
\end{scope}
\end{tikzpicture}
\end{center}

\emph{同一个集合，粗粒度下切开了信息块因而无法询问，细粒度下恰是若干块的并，才成为事件。}

换一台能分辨全部四个结果的设备，可回答的事件族扩张成幂集 \(2^\Omega\)。包含关系
\[
\mathcal G\subseteq\mathcal F
\]
读作``\(\mathcal F\) 至少和 \(\mathcal G\) 一样精细''。随时间累积信息的 filtration
\[
\mathcal F_0\subseteq\mathcal F_1\subseteq\cdots
\]
把``知道得越来越多''写成了集合族的递增。日志加一个字段、埋点多打一个维度，做的都是同一件事：把图右边那些竖线添进去。

\(\sigma\)-代数在连续空间里同时承担两个角色：划定哪些集合有定义完备的概率，以及标记当前信息能分辨到什么程度。第二个角色到第五节讲条件期望时会全部兑现。

\subsection{可数并是极限操作的最低价格}

\begin{definition}
给定样本空间 \(\Omega\)，集合族 \(\mathcal F\subseteq 2^\Omega\) 称为 \(\Omega\) 上的 \textbf{\(\sigma\)-代数}，如果它满足：\(\Omega\in\mathcal F\)；若 \(A\in\mathcal F\) 则 \(A^{c}\in\mathcal F\)；若 \(A_1,A_2,\ldots\in\mathcal F\) 则 \(\bigcup_{n\ge 1}A_n\in\mathcal F\)。
\end{definition}

三条规则合起来说的是一件事：凡是用可数多次``且、或、非''从已有事件拼出来的新事件，仍然留在族内。你不会拼着拼着掉出可询问的范围。

第三条里的关键词是可数，不是有限。有限并看着已经够用，直到极限出现。概率论中反复出现的事件形如
\[
\bigcup_{n=1}^{\infty}A_n,\qquad
\bigcap_{n=1}^{\infty}A_n,\qquad
\limsup_{n\to\infty} A_n=\bigcap_{m\ge1}\bigcup_{n\ge m}A_n,
\]
最后那个读作``\(A_n\) 发生无穷多次''。Borel--Cantelli、几乎必然收敛、首次击中时间，这些论证的对象全是这类可数拼装出来的集合。只对有限并封闭的集合族，等于宣布这些问题不许问，而它们恰是概率论最想问的问题。

另一头也有界限：要求对任意不可数并封闭就太强了。实数轴上每个单点集都是 Borel 集，若不可数并也封闭，那么任何子集都能写成它含有的单点的并，\(\mathcal B(\R)\) 就退化成幂集，Vitali 的障碍立刻找上门。可数是极限论证够用、而可加性还撑得住的那个刻度。

最粗的 \(\sigma\)-代数是 \(\set{\varnothing,\Omega}\)，只区分必然与不可能，对应一份什么都没记的日志。最细的是幂集。连续空间里用的那个落在两者之间。

\subsection{生成运算让你只指定少数几个集合}

没人手工写出可数次补与并之后的全部结果。做法是指定一族起始集合 \(\mathcal C\)，取所有包含 \(\mathcal C\) 的 \(\sigma\)-代数的交，得到包含 \(\mathcal C\) 的最小 \(\sigma\)-代数，记作 \(\sigma(\mathcal C)\)。任意多个 \(\sigma\)-代数的交仍是 \(\sigma\)-代数，所以这个交确实存在，且幂集保证了候选非空。这套路子与线性代数里``一组向量张成的子空间''完全同构：都是取包含给定对象的最小封闭结构。

实轴上最重要的一例由全体开区间生成，
\[
\mathcal B(\R)=\sigma\bigl(\set{(a,b):a<b}\bigr),
\]
称为 Borel \(\sigma\)-代数。它装下了开集、闭集、可数集，以及从开区间出发经可数次集合运算能到达的一切，但装不下 \(\R\) 的全部子集。连续分布与实值随机变量都在 \(\mathcal B(\R)\) 上定义，见 \hyperref[sec:05-Probability/03-Random-Variables/03]{``连续分布、密度与分位数''}。

\subsection{可测函数不制造新的不可回答问题}

从 \((\Omega,\mathcal F)\) 到 \((\R,\mathcal B(\R))\) 的函数 \(X\)，若对每个 \(B\in\mathcal B(\R)\) 都有
\[
X^{-1}(B)=\set{\omega\in\Omega : X(\omega)\in B}\in\mathcal F,
\]
就称为可测函数。检验时不必遍历所有 Borel 集：只要验证所有半无限区间 \((-\infty,x]\) 的原像在 \(\mathcal F\) 内即可。这些区间生成 \(\mathcal B(\R)\)，而原像运算与补、可数并都交换，所以可测性顺着生成关系一路传上去。

随机变量必须可测，理由在上一段的形式里已经写着了。问``\(X>3\) 的概率是多少''，实际是在问集合 \(\set{\omega:X(\omega)>3}\) 的概率；这个集合若不在 \(\mathcal F\) 中，\(P\) 压根没定义在它上面，问题悬空。可测性保证了：关于 \(X\) 取值的任何 Borel 问题，都能拉回成原空间里有概率的事件。

可测性比连续性弱得多。连续函数一定 Borel 可测，但大量处处不连续的函数同样可测，比如有理数集的指示函数。可测只要求原像保住可询问性，不要求保开集或保闭集。

随机变量自己携带的信息记作
\[
\sigma(X)=\set{X^{-1}(B) : B\in\mathcal B(\R)},
\]
它恰好由``只看 \(X\) 的取值就能判定''的那些事件组成。若 \(Y=g(X)\) 且 \(g\) 可测，则
\[
\sigma(Y)\subseteq\sigma(X),
\]
因为对 \(X\) 做确定性变换产生不了新信息：关于 \(Y\) 的任何问题都可以翻译成关于 \(X\) 的问题。这条包含关系的方向，正是判断一次特征工程、一次降维、一个统计量是否丢掉了要紧信息的入口——包含是必然的，问题只在于丢了多少。

\subsection{不可测与概率为零是两码事}

事件不属于 \(\mathcal F\)，说的是模型没为它定义概率。这与概率已定义、恰好等于零，完全不同。

零概率事件是合法的空间成员，只是被赋了零测度；它甚至可以含有不可数个样本点，比如 \([0,1]\) 上均匀分布下的有理数集，或者任何一个单点。不可测集连进入模型的资格都没有。前者是``问了，答案是 0''，后者是``这个框架下问不出口''。

区别在信息变动时更明显。filtration 扩张，原先不可测的集合可能变得可测——回到本节开头那张图，添一条竖线，灰色区域就成了事件。零概率却不会因为信息增多而变正：测度为零是测度的属性，与你观测得多细无关。日志加字段能让你问出更多问题，但不会把一件不可能的事变成可能。

连续概率模型的建造顺序因此是固定的：先定 \(\sigma\)-代数，声明哪些集合可测；再在这些集合上定义测度。跳过第一步直接谈概率，等于在 Vitali 集上要求长度，两个条件在动笔之前就已经打架了。下一节从第二步开始：一个可数可加的非负集函数，怎样把长度、计数与概率压成同一个对象。
